\chapter{Gaussian states and the covariance description}
\label{narr:gaussian}

Gaussian states replace exponentially many amplitudes by a two-point matrix.
This economy is powerful because higher correlations then follow by a specified
factorization. It is also the approximation's boundary: the factorization is a
property of the chosen state, not an operator identity obeyed by every spin
state or every projected fermion state.

% EDITORIAL QUESTION: Which signs and occupations does our covariance convention encode?
\section{Fixing the convention with one complex mode}
For $2n$ Majoranas define
\begin{equation}
 \Omega_{ab}=\frac{\ii}{2}\Tr\rho[\gamma_a,\gamma_b],
 \qquad \langle\gamma_a\gamma_b\rangle=\delta_{ab}-\ii\Omega_{ab}.
 \label{eq:narr:covariance}
\end{equation}
Then $\Omega$ is real and antisymmetric. For one mode
$f=(\gamma_1+\ii\gamma_2)/2$, write
$J_2=\left(\begin{smallmatrix}0&1\\-1&0\end{smallmatrix}\right)$.
If $p=\langle f^\dagger f\rangle$, then
\begin{equation}
 \Omega=(2p-1)J_2,\qquad
 H=\frac{\ii}{4}\bm\gamma^T(\varepsilon J_2)\bm\gamma
   =\varepsilon(f^\dagger f-\tfrac12).
 \label{eq:narr:one_mode}
\end{equation}
For $\varepsilon>0$ the empty mode is the ground state and has $\Omega=-J_2$.
This small example fixes the thermal sign and the factor between matrix bands
and occupation energies without referring to a dispersion plot.

% EDITORIAL QUESTION: Which two-point matrices describe Gaussian states, and when are they pure?
\section{Positivity, purity, and the state manifold}
Here Gaussian means a parity-even quasifree state, including singular pure
limits, determined by its covariance and Wick rule. Orthogonal block reduction
gives $\Omega=O\bigoplus_j\lambda_jJ_2\,O^T$. Such a Gaussian density operator
exists precisely when $|\lambda_j|\leq1$; its independent-mode probabilities
are $(1+\lambda_j)/2$. It is pure precisely when every probability is zero or
one, equivalently
\begin{equation}
 \Omega^2=-\Id.
 \label{eq:narr:purity}
\end{equation}
The pure covariances form the space of orthogonal complex structures
\begin{equation}
 \mathcal J_{2n}=\{\Omega\in\RR^{2n\times2n}:\Omega^T=-\Omega,
       \ \Omega^2=-\Id\}\simeq O(2n)/U(n).
 \label{eq:narr:complex_structures}
\end{equation}
Its real dimension is $n(n-1)$, and it has two connected components. This is
not the ordinary real Grassmannian of $n$-planes in $\RR^{2n}$, whose dimension
is $n^2$. Older records used the same name for these different objects; the
explicit defining equations avoid that ambiguity.

For a pure covariance, $\Pi_{\rm occ}=(\Id+\ii\Omega)/2$ is a rank-$n$
\emph{one-particle} projector. It is not a many-body ground-space or flux-sector
projector. In a Fourier representation the reality of real-space Majoranas
links $\bk$ and $-\bk$. A decomposition into momentum blocks must retain that
relation, including special momenta fixed by inversion; one cannot simply
assign an independent real covariance manifold to every generic momentum.

% EDITORIAL QUESTION: What does Wick factorization compute and what can it miss?
\section{Ordered correlations and Pfaffians}
For distinct indices in the stated order, the Gaussian Wick rule reads
\begin{equation}
 \langle\gamma_{a_1}\cdots\gamma_{a_{2m}}\rangle
       =(-\ii)^m\Pf\Omega_{\{a_1,\ldots,a_{2m}\}}.
 \label{eq:narr:wick}
\end{equation}
For four indices this is
$-(\Omega_{ab}\Omega_{cd}-\Omega_{ac}\Omega_{bd}+\Omega_{ad}\Omega_{bc})$.
Permuting the Majoranas changes the sign according to the CAR. For repeated
indices first reduce the product with the CAR; Eq.~\eqref{eq:narr:wick} was
stated for distinct ones.

This rule allows a spin correlator represented by an even Majorana product to
be evaluated in an unprojected Gaussian state. It does not justify replacing
an exact expectation by a product of link expectations, nor does it turn an
auxiliary flux into a conserved physical sector. For the canonical ordered
parity operator $\mathcal P=(-\ii)^n\gamma_1\cdots\gamma_{2n}$,
$\langle\mathcal P\rangle=(-1)^n\Pf\Omega$. Ordering is part of this formula.
The projection example in Chapter~\ref{narr:physical} shows why a covariance
alone need not determine a physical projected four-point function.
